\documentclass{article}
\usepackage{xcolor}
\usepackage{mdframed}
\usepackage{propositions}

%% A "boxed" style: wraps the whole prop list in an mdframed environment.
%% The wrapper begin/end keys are environment-level: they fire once around
%% \begin{list}...\end{list}, not once per item.  Values are braced so any
%% commas inside the mdframed options are not seen as key separators.
\SetPropStyle{boxed}{
  wrapper begin = {\begin{mdframed}},
  wrapper end   = {\end{mdframed}},
}

%% A coloured variant, to show options can be passed to the wrapper env.
\SetPropStyle{shaded}{
  wrapper begin = {\begin{mdframed}[backgroundcolor=blue!8, linecolor=blue!40]},
  wrapper end   = {\end{mdframed}},
}

%% Inheritance: a style whose parent is "boxed" inherits the wrapper and
%% adds item-level formatting (bold names).  Demonstrates that wrapper keys
%% propagate down the style chain.
\SetPropStyle{boxedbold}{
  style  = boxed,
  format = \textbf{#1},
}

\begin{document}

\section{Boxed style via \texttt{style=boxed}}

The entire list below should sit inside a single framed box:

\begin{prop}[style=boxed]
  \item First proposition. \label{w:first}
  \item Second proposition. \label{w:second}
  \item[Key Claim] A named item inside the box. \label{w:named}
\end{prop}

Cross-references still work across the box boundary:
\ref{w:first}, \ref{w:second}, \ref{w:named}.

\section{Coloured wrapper with options}

\begin{prop}[style=shaded]
  \item This list is in a shaded, blue-edged box.
  \item Second item.
\end{prop}

\section{Inheriting a wrapper (\texttt{boxedbold})}

\begin{prop}[style=boxedbold]
  \item[Lemma] Inherits the box from \texttt{boxed}, names are bold.
  \item[Corollary] Another bold-named, boxed item.
\end{prop}

\section{Wrapper keys used directly (no style)}

\begin{prop}[wrapper begin = {\begin{mdframed}}, wrapper end = {\end{mdframed}}]
  \item Boxed via direct \texttt{wrapper begin}/\texttt{wrapper end} keys.
  \item Second item.
\end{prop}

\section{An ordinary list is not boxed}

This list should have no frame (the env-arg wrapper above was local to its
own environment):

\begin{prop}
  \item Plain, unboxed proposition.
  \item Another plain one.
\end{prop}

\section{Nesting: only the outermost list is boxed}

Only the outer list should get a frame; the sublist sits inside that same
frame and is not separately boxed.

\begin{prop}[style=boxed]
  \item Outer item, inside the box. \label{w:outer}
  \begin{prop}
    \item Inner item (sublist is not separately boxed).
    \item Second inner item.
  \end{prop}
  \item Back to the outer level. \label{w:back}
\end{prop}

References: \ref{w:outer}, \ref{w:back}.

A standard \texttt{itemize} nested in a boxed prop keeps its bullets and is
not boxed:

\begin{prop}[style=boxed]
  \item A proposition.
  \begin{itemize}
    \item an ordinary bullet
    \item another bullet
  \end{itemize}
  \item Another proposition.
\end{prop}

\section{\texttt{\textbackslash propoptions\{style=boxed\}} as a scoped default}

\begingroup
\propoptions{style=boxed}

Both lists in this group are boxed by default:

\begin{prop}
  \item Boxed by the group-level default.
\end{prop}

\begin{prop}
  \item Also boxed by the same default.
\end{prop}
\endgroup

Outside the group, the default is gone, so this list is unboxed:

\begin{prop}
  \item Not boxed.
\end{prop}

\section{Nesting wrappers explicitly}

A sublist that requests a wrapper of its own is boxed separately, inside the
outer box.  A \emph{different} wrapper (plain frame outside, shaded inside):

\begin{prop}[style=boxed]
  \item Outer item. \label{w:no}
  \begin{prop}[style=shaded]
    \item Inner, shaded.
    \item Inner, shaded two.
  \end{prop}
  \item Outer item again.
\end{prop}

The \emph{same} wrapper explicitly re-requested also nests (frame inside a
frame); only an \emph{inherited}, unrequested wrapper is suppressed:

\begin{prop}[style=boxed]
  \item Outer item.
  \begin{prop}[style=boxed]
    \item Inner, its own frame.
  \end{prop}
  \item Outer item again.
\end{prop}

\section{\texttt{\textbackslash propoptions} inside a boxed environment}

\texttt{\textbackslash propoptions\{style=shaded\}} issued inside the outer
(boxed) environment---no explicit group needed---makes each following sublist
shaded; the outer box is unaffected and still closes correctly:

\begin{prop}[style=boxed]
  \item Outer, before switching the sublist default.
  \propoptions{style=shaded}
  \begin{prop}
    \item First shaded sublist.
  \end{prop}
  \begin{prop}
    \item Second shaded sublist.
  \end{prop}
  \item Outer, after (still only the outer box).
\end{prop}

\end{document}
